\documentclass[11pt]{article}
\usepackage[T1]{fontenc}
\usepackage{lmodern}
\usepackage[a4paper,margin=25mm]{geometry}
\usepackage{amsmath,amssymb,amsthm}
\usepackage[hypertexnames=false,hidelinks]{hyperref}
\newcommand{\V}{\mathsf V}
\newcommand{\E}{\mathsf E}
\newcommand{\PE}{\Phi_{\E}}
\newcommand{\PV}{\Phi_{\V}}
\newcommand{\flow}{\mathsf{flow}_{+}}
\newcommand{\loc}[1]{\par\smallskip\noindent{\small\textit{Source: #1}}\par\smallskip}
\setlength{\emergencystretch}{2em}
\begin{document}
\begin{center}
{\large Probability on Trees and Networks\\[2pt]Exercise 6.27 and its Author Comment}\\[4pt]
Russell Lyons\quad Yuval Peres\\[2pt]
Version of 19 August 2026---corrected and updated paperback edition
\end{center}

% BEGIN CONTEXT 6.27
\noindent\textbf{Source and scope (editorial).}
This private reader reproduces selected English author text from the fixed 716-page paperback PDF (21,010,205 bytes). Exercise 6.27 is on physical p.\ 245 (printed p.\ 228); its semantic target is the complete statement of Theorem 6.2 on physical p.\ 196 (printed p.\ 179). Its complete Comment is on physical p.\ 633 (printed p.\ 616). The supporting Theorem 6.1 and its complete printed global-flow construction are reprinted below as uncounted context. The body theorem and the exercise are one target. The refinement in Exercise 6.28 is excluded.

The book authors are Russell Lyons and Yuval Peres. The Comment expressly identifies its route as the original proof in BLPS (1999b); this is a historical attribution within the book, not a claim to have acquired or transcribed the separate paper. The original book and the separate errata dated 19 August 2026 remain distinct source materials. The erratum for printed p.\ 179 concerns the unselected submodularity proof of Theorem 6.2; this reader supplies the prescribed flow proof instead. No source mathematics is repaired or independently certified.

\medskip\noindent\textbf{Oriented-edge notation.}
\loc{Section 1.1, physical p.\ 19 (printed p.\ 2); Section 2.4, physical p.\ 49 (printed p.\ 32).}
The endpoints $e^-$ and $e^+$ are the tail and head of $e$. The edge $-e$ has the reverse orientation. Each edge occurs with both orientations; an antisymmetric function satisfies $\theta(-e)=-\theta(e)$. The book's boundary-operator formula for net flow out of a vertex is
\[
(d^*\theta)(x):=\sum_{e^-=x}\theta(e).
\]
These are notation excerpts; no square-summability restriction is added to the antisymmetric functions quantified in Theorem 6.1 or the Comment.

\subsection*{Network and expansion definitions}
\loc{Physical p.\ 194 (printed p.\ 177).}
Let $G$ be a connected graph. Weight the edges by positive numbers $c(e)$ and the vertices by positive numbers $D(x)$. We call $(G,c,D)$ a network, which is a little more general than the networks we have used before where we had weights only for the edges. We assume that $G$ is locally finite, or, more generally, that for each $x$,
\begin{equation}
\sum_{e^-=x}c(e)<\infty.\tag{6.1}
\end{equation}
For $K\subset\V$, we define the edge boundary $\partial_{\E}K$ to be the set of (unoriented) edges that connect $K$ to its complement. Write
\[
|K|_D:=\sum_{x\in K}D(x),\qquad
|\partial_{\E}K|_c:=\sum_{e\in\partial_{\E}K}c(e).
\]
Define the edge-expansion constant or edge-isoperimetric constant of $(G,c,D)$ by
\[
\PE(G):=\PE(G,c,D):=
\inf\left\{\frac{|\partial_{\E}K|_c}{|K|_D};\ \varnothing\ne K\subset\V\text{ is finite}\right\}.
\]
\loc{Transitivity definition, physical p.\ 196 (printed p.\ 179).}
Here, we say that a network $(G,c,D)$ is transitive if, for every pair $x,y\in\V(G)$, there is an automorphism of the graph $G$ that takes $x$ to $y$ and that preserves the edge weights $c$ and vertex weights $D$.

\subsection*{The complete supporting flow construction}
\loc{Theorem 6.1 and its proof, physical p.\ 195 (printed p.\ 178).}
The following theorem (and Exercise 6.2) shows that $\PE(G)$ is precisely the maximum proportional profit $\alpha$ that we seek (actually, because of our choice of $d^*$, we need to take the negative of the function $\theta$ guaranteed by Theorem 6.1). The theorem thus dualizes the infimum in the definition of $\PE(G)$ not only to a supremum, but to a maximum. It is due to Benjamini, Lyons, Peres, and Schramm (1999b), hereinafter referred to as BLPS (1999b).

\medskip\noindent\textbf{Theorem 6.1 (Duality for Edge Expansion).}
\textit{For any network $(G,c,D)$, we have}
\[
\begin{split}
\PE(G,c,D)=\max\bigl\{\alpha\geq0;\ &\exists\theta\ \forall e\ |\theta(e)|\leq c(e)\text{ and }\\
&\forall x\ d^*\theta(x)=\alpha D(x)\bigr\},
\end{split}
\]
\textit{where $\theta$ runs over the antisymmetric functions on $\E$.}
\begin{proof}
Denote by $A$ the set of $\alpha$ of which the maximum is taken on the right-hand side of the desired equation. Thus, we want to show that $\max A=\PE(G)$. In fact, we will show that $A=[0,\PE(G)]$ as a consequence of the max-flow min-cut theorem.

Given $\alpha\geq0$ and any finite, nonempty $K\subset\V$, define the network $G(K)=G(K,\alpha)$ with vertices $K$ and two extra vertices, $a$ and $z$. The edges are those of $G$ with both endpoints in $K$, those in $\partial_{\E}K$ where the endpoint not in $K$ is replaced by $z$, and an edge between $a$ and each point in $K$. Give each edge $e$ with both endpoints in $K$ the capacity $c(e)$. Give each edge $[x,z]$ the capacity of the corresponding edge in $\partial_{\E}K$. Give the other edges capacity $c(a,x):=\alpha D(x)$.

In view of the cutset consisting of all edges incident to $z$, it is clear that any admissible flow from $a$ to $z$ in $G(K)$ has strength at most $|\partial_{\E}K|_c$. Suppose that $\alpha\in A$, in other words, that there is a function $\theta$ satisfying $\forall e\ |\theta(e)|\leq c(e)$ and $\forall x\ d^*\theta(x)=\alpha D(x)$. This function $\theta$ induces an admissible flow on $G(K)$ from $a$ to $z$ of strength $\alpha|K|_D$ by putting the flow $\alpha D(x)$ on each edge $(a,x)$ for $x\in K$, whence $\alpha\leq|\partial_{\E}K|_c/|K|_D$. Since this holds for all $K$, we get $\alpha\leq\PE(G)$.

In the other direction, if $\alpha\leq\PE(G)$, then for all finite $K$, we claim that there is an admissible flow from $a$ to $z$ in $G(K)$ of strength $\alpha|K|_D$. Consider any cutset $\Pi$ separating $a$ and $z$. Let $K'$ be the vertices in $K$ that $\Pi$ separates from $z$. Then $\partial_{\E}K'\subseteq\Pi$ and $[a,x]\in\Pi$ for all $x\in K\setminus K'$. Therefore,
\[
\begin{aligned}
\sum_{e\in\Pi}c(e)
&\geq|\partial_{\E}K'|_c+\alpha|K\setminus K'|_D\\
&\geq\PE(G)|K'|_D+\alpha|K\setminus K'|_D\\
&\geq\alpha|K'|_D+\alpha|K\setminus K'|_D=\alpha|K|_D.
\end{aligned}
\]
Thus, our claim follows from the max-flow min-cut theorem. Note that the flow along $[a,x]$ is $\alpha D(x)$ for every $x\in K$, since the flow's strength is $\alpha|K|_D$.

Now let $K_n$ be finite sets increasing to $\V$, and let $\theta_n$ be admissible flows on $G(K_n)$ with strength $\alpha|K_n|_D$. There is a subsequence $\langle n_i\rangle$ such that for all edges $e\in\E$, the limit $\theta(e):=\theta_{n_i}(e)$ exists; clearly $\forall e\ |\theta(e)|\leq c(e)$ and, by (6.1) and the dominated convergence theorem, $\forall x\ d^*\theta(x)=\alpha D(x)$. Thus, $\alpha\in A$.
\end{proof}
% END CONTEXT 6.27

\clearpage
% BEGIN PRIMARY STATEMENT 6.27
\section*{Exercise 6.27 and its full target}
\loc{Exercise 6.27: physical p.\ 245 (printed p.\ 228).}
\noindent\textbf{Exercise 6.27.} Use Theorem 6.1 and its proof to give another proof of Theorem 6.2.

\loc{Theorem 6.2: physical p.\ 196 (printed p.\ 179).}
\noindent\textbf{Theorem 6.2.}
\textit{If $(G,c,D)$ is an infinite, transitive network, then for all finite, nonempty $K\subset\V$, we have $|\partial_{\E}K|_c/|K|_D>\PE(G)$.}
% END PRIMARY STATEMENT 6.27

% BEGIN PRIMARY PROOF 6.27
\section*{Complete author Comment on Exercise 6.27}
\loc{Physical p.\ 633 (printed p.\ 616).}
(This was the original proof of Theorem 6.2 in BLPS (1999b).) Let $\theta$ be an antisymmetric function on $\E$ with $|\theta(e)|\leq c(e)$ for all edges $e$ and $d^*\theta(x)=\PE(G)D(x)$ for all vertices $x$. We may assume that $\theta$ is acyclic, that is, that there is no cycle of oriented edges on each of which $\theta>0$. (Otherwise, we modify $\theta$ by subtracting appropriate cycles.) Let $K\subset\V$ be finite and nonempty. Suppose for a contradiction that $|\partial_{\E}K|_c/|K|_D=\PE(G)$. The proof of Theorem 6.1 shows that for all $e\in\partial_{\E}K$, we have $\theta(e)=c(e)$ if $e$ is oriented to point out of $K$; in particular, $\theta(e)>0$. Let $(x_1,x_0)\in\partial_{\E}K$ with $x_1\in K$ and $x_0\notin K$, and let $\gamma$ be an automorphism of $G$ that carries $x_0$ to $x_1$. Write $x_2$ for the image of $x_1$ and $\gamma K$ for the image of $K$. Since we also have $|\partial_{\E}\gamma K|_c/|\gamma K|_D=\PE(G)$ and $(x_2,x_1)\in\partial_{\E}\gamma K$, it follows that $\theta(x_2,x_1)>0$. We may similarly find $x_3$ such that $\theta(x_3,x_2)>0$ and so on, until we arrive at some $x_k$ that equals some previous $x_j$ or is outside $K$. Both lead to a contradiction, the former contradicting the acyclicity of $\theta$ and the latter the fact that on all edges leading out of $K$, we have $\theta>0$.
% END PRIMARY PROOF 6.27

\medskip\noindent\textbf{Cached bibliographic attribution (not an additional source proof).}
The book bibliography, physical p.\ 667 (printed p.\ 650), lists: Benjamini, I., Lyons, R., Peres, Y., and Schramm, O.\ (1999b), \textit{Group-invariant percolation on graphs}, \textit{Geom.\ Funct.\ Anal.}\ \textbf{9}(1), 29--66, MR: 99m:60149.
\end{document}
